\documentclass[11pt]{article}
\newcommand{\ListaTitulo}{Lista 14: Simulado Prova 3 (Unidade 4)}
% Preâmbulo compartilhado: MATF14, Listas de Exercícios 2026
% Usado por ListaNN.tex e GabaritoNN.tex via % Preâmbulo compartilhado: MATF14, Listas de Exercícios 2026
% Usado por ListaNN.tex e GabaritoNN.tex via % Preâmbulo compartilhado: MATF14, Listas de Exercícios 2026
% Usado por ListaNN.tex e GabaritoNN.tex via \input{preamble}

\usepackage[utf8]{inputenc}
\usepackage[T1]{fontenc}
\usepackage[brazil]{babel}
\usepackage{fullpage}
\usepackage{amsmath,amssymb}
\usepackage{enumitem}
\usepackage{xcolor}
\usepackage{fancyhdr}
\usepackage{titlesec}
\usepackage{listings}
\usepackage{hyperref}
\usepackage{booktabs}
\usepackage{graphicx}
\usepackage{tikz}

\definecolor{matf14azul}{HTML}{0D3B54}
\definecolor{matf14dourado}{HTML}{C9971E}

\titleformat{\section}{\normalfont\Large\bfseries\color{matf14azul}}{\thesection}{1em}{}
\titleformat{\subsection}{\normalfont\large\bfseries\color{matf14azul}}{\thesubsection}{1em}{}

\providecommand{\ListaTitulo}{Lista de Exercícios}

\setlength{\headheight}{14pt}
\pagestyle{fancy}
\fancyhf{}
\lhead{\small MATF14: Estatística Econômica I}
\rhead{\small \ListaTitulo}
\cfoot{\thepage}
\renewcommand{\headrulewidth}{0.4pt}

\lstset{
  basicstyle=\ttfamily\small,
  breaklines=true,
  frame=single,
  columns=fullflexible,
  backgroundcolor=\color{gray!8},
  commentstyle=\color{matf14dourado},
  keywordstyle=\color{matf14azul}\bfseries,
  language=R
}

\newcounter{questaocounter}
\newenvironment{questao}[1]{%
  \stepcounter{questaocounter}%
  \bigskip\noindent\textbf{Questão #1.}\ %
}{\par}

\newcommand{\cabecalho}[2]{%
  \begin{center}
    {\Large\bfseries\color{matf14azul} #1}\\[0.3em]
    {\large #2}
  \end{center}
  \vspace{0.5em}
  \noindent\rule{\linewidth}{0.6pt}
  \vspace{0.5em}
}

\newenvironment{solucao}{%
  \par\smallskip\noindent\textit{\color{matf14dourado!80!black}Solução.}\ %
}{\hfill$\blacksquare$\par}


\usepackage[utf8]{inputenc}
\usepackage[T1]{fontenc}
\usepackage[brazil]{babel}
\usepackage{fullpage}
\usepackage{amsmath,amssymb}
\usepackage{enumitem}
\usepackage{xcolor}
\usepackage{fancyhdr}
\usepackage{titlesec}
\usepackage{listings}
\usepackage{hyperref}
\usepackage{booktabs}
\usepackage{graphicx}
\usepackage{tikz}

\definecolor{matf14azul}{HTML}{0D3B54}
\definecolor{matf14dourado}{HTML}{C9971E}

\titleformat{\section}{\normalfont\Large\bfseries\color{matf14azul}}{\thesection}{1em}{}
\titleformat{\subsection}{\normalfont\large\bfseries\color{matf14azul}}{\thesubsection}{1em}{}

\providecommand{\ListaTitulo}{Lista de Exercícios}

\setlength{\headheight}{14pt}
\pagestyle{fancy}
\fancyhf{}
\lhead{\small MATF14: Estatística Econômica I}
\rhead{\small \ListaTitulo}
\cfoot{\thepage}
\renewcommand{\headrulewidth}{0.4pt}

\lstset{
  basicstyle=\ttfamily\small,
  breaklines=true,
  frame=single,
  columns=fullflexible,
  backgroundcolor=\color{gray!8},
  commentstyle=\color{matf14dourado},
  keywordstyle=\color{matf14azul}\bfseries,
  language=R
}

\newcounter{questaocounter}
\newenvironment{questao}[1]{%
  \stepcounter{questaocounter}%
  \bigskip\noindent\textbf{Questão #1.}\ %
}{\par}

\newcommand{\cabecalho}[2]{%
  \begin{center}
    {\Large\bfseries\color{matf14azul} #1}\\[0.3em]
    {\large #2}
  \end{center}
  \vspace{0.5em}
  \noindent\rule{\linewidth}{0.6pt}
  \vspace{0.5em}
}

\newenvironment{solucao}{%
  \par\smallskip\noindent\textit{\color{matf14dourado!80!black}Solução.}\ %
}{\hfill$\blacksquare$\par}


\usepackage[utf8]{inputenc}
\usepackage[T1]{fontenc}
\usepackage[brazil]{babel}
\usepackage{fullpage}
\usepackage{amsmath,amssymb}
\usepackage{enumitem}
\usepackage{xcolor}
\usepackage{fancyhdr}
\usepackage{titlesec}
\usepackage{listings}
\usepackage{hyperref}
\usepackage{booktabs}
\usepackage{graphicx}
\usepackage{tikz}

\definecolor{matf14azul}{HTML}{0D3B54}
\definecolor{matf14dourado}{HTML}{C9971E}

\titleformat{\section}{\normalfont\Large\bfseries\color{matf14azul}}{\thesection}{1em}{}
\titleformat{\subsection}{\normalfont\large\bfseries\color{matf14azul}}{\thesubsection}{1em}{}

\providecommand{\ListaTitulo}{Lista de Exercícios}

\setlength{\headheight}{14pt}
\pagestyle{fancy}
\fancyhf{}
\lhead{\small MATF14: Estatística Econômica I}
\rhead{\small \ListaTitulo}
\cfoot{\thepage}
\renewcommand{\headrulewidth}{0.4pt}

\lstset{
  basicstyle=\ttfamily\small,
  breaklines=true,
  frame=single,
  columns=fullflexible,
  backgroundcolor=\color{gray!8},
  commentstyle=\color{matf14dourado},
  keywordstyle=\color{matf14azul}\bfseries,
  language=R
}

\newcounter{questaocounter}
\newenvironment{questao}[1]{%
  \stepcounter{questaocounter}%
  \bigskip\noindent\textbf{Questão #1.}\ %
}{\par}

\newcommand{\cabecalho}[2]{%
  \begin{center}
    {\Large\bfseries\color{matf14azul} #1}\\[0.3em]
    {\large #2}
  \end{center}
  \vspace{0.5em}
  \noindent\rule{\linewidth}{0.6pt}
  \vspace{0.5em}
}

\newenvironment{solucao}{%
  \par\smallskip\noindent\textit{\color{matf14dourado!80!black}Solução.}\ %
}{\hfill$\blacksquare$\par}


\begin{document}

\cabecalho{Lista de Exercícios 14: Simulado}{Revisão geral, no formato e no nível da Prova 3: Variáveis Aleatórias Contínuas (Cap. 4, Aulas 24--30)}

\noindent \textit{Simulado de 10 pontos, no mesmo formato da prova oficial: questões com múltiplos
itens, valor indicado entre parênteses. Resolva com tempo cronometrado (sugestão: 100 minutos,
sem consulta) e traga as dúvidas para a aula de revisão.}

\begin{enumerate}

\item Seja $C>0$ uma constante e $X$ uma variável aleatória contínua com densidade
$$
f(x) = \begin{cases} Cx^{-2}, & x\ge 5 \\ 0, & x<5 \end{cases}
$$
\begin{enumerate}
  \item (1 ponto) Encontre o valor de $C$ para que $f$ seja uma densidade de probabilidade válida.
  \item (1 ponto) Calcule a função de distribuição acumulada $F(x)$, para $x\ge5$.
  \item (1 ponto) Calcule $P(X>10)$.
\end{enumerate}

\item O tempo de espera (em minutos) para atendimento em uma agência bancária é
$X\sim\text{Unif}(0,20)$.
\begin{enumerate}
  \item (1 ponto) Calcule $E(X)$ e $\mathrm{Var}(X)$.
  \item (1 ponto) Calcule a probabilidade de um cliente esperar entre 5 e 12 minutos.
\end{enumerate}

\item O retorno anual de um fundo multimercado é aproximadamente Normal, com média $10\%$ e
desvio padrão $18\%$.
\begin{enumerate}
  \item (1 ponto) Calcule a probabilidade do retorno em um ano ser maior que $15\%$.
  \item (1 ponto) Calcule a probabilidade do retorno em um ano ser menor que $-8\%$.
\end{enumerate}

\item O peso (em kg) de pacotes de um produto tem densidade
$$
f(x) = \begin{cases} \frac{3}{8}x^2, & 0\le x\le 2 \\ 0, & \text{caso contrário} \end{cases}
$$
\begin{enumerate}
  \item (1 ponto) Confirme que $f$ é uma densidade válida.
  \item (1 ponto) Calcule $E(X)$.
  \item (1 ponto) Calcule $P(X\le 1)$.
\end{enumerate}

\end{enumerate}

\vspace{1cm}
\begin{center}
\Large\textbf{Formulário}
\end{center}
\begin{small}
\begin{itemize}
  \item $P(a\le X\le b)=\int_a^b f(x)dx$; $\int_{-\infty}^{\infty} f(x)dx=1$.
  \item $E(X)=\int xf(x)dx$; $\mathrm{Var}(X)=E(X^2)-[E(X)]^2$; $dp(X)=\sqrt{\mathrm{Var}(X)}$.
  \item $F(x)=\int_{-\infty}^x f(t)dt$; $F'(x)=f(x)$.
  \item $X\sim\text{Unif}(\alpha,\beta)$: $f(x)=1/(\beta-\alpha)$; $E(X)=(\alpha+\beta)/2$,
    $\mathrm{Var}(X)=(\beta-\alpha)^2/12$.
  \item $X\sim N(\mu,\sigma^2)$: $E(X)=\mu$, $\mathrm{Var}(X)=\sigma^2$; $Z=(X-\mu)/\sigma\sim
    N(0,1)$.
  \item $\int x^n\,dx = \dfrac{x^{n+1}}{n+1}+c$ ($n\ne-1$); $\displaystyle\int e^{ax}\,dx =
    \dfrac1a e^{ax}+c$.
\end{itemize}
\end{small}


\bigskip
\section*{Exercícios complementares}

\begin{enumerate}[label=\textbf{C\arabic*.}, leftmargin=*]
\item (Prova 2 de 2026.1) Seja $f(x) = Cx^{-3}$ para $x \geq 10$ e zero caso contrário.
\begin{enumerate}[label=(\alph*)]
  \item Encontre $C$.
  \item Calcule a função de distribuição acumulada de $X$.
  \item Calcule $P(X > 20)$.
\end{enumerate}
\item (Prova 2 de 2026.1) O retorno anual de um fundo é normal com média 8\% e desvio padrão 20\%. Calcule a probabilidade de o retorno ser maior que 10\% e de ser menor que $-5\%$.
\item Considere $f(x) = 3x^2$ para $-1 \leq x \leq 0$ e zero caso contrário. Calcule $E(X)$ e $Var(X)$.

\end{enumerate}

\end{document}
